\documentclass{article}

\usepackage{amsmath,amssymb}
\usepackage{xurl}
\usepackage[hidelinks]{hyperref}
\setlength{\emergencystretch}{2em}

\input{command}

\title{Wolf Equation~(8.103) --- Small-Power Exponent Scope}
\author{}
\date{2026-08-09}

\begin{document}
\maketitle
\gapnote{scope-restriction}{resolved}

\section{The source statement}

Wolf considers an upper-triangular matrix $X=\Lambda+N\in M_D(\mathbb C)$,
where $\Lambda$ is diagonal and $N$ is strictly upper triangular.  For an
arbitrary submultiplicative norm satisfying $\lVert\Lambda\rVert\le 1$, the
notes state for every $n\in\mathbb N$ that
\begin{align}
  \lVert X^n\rVert
  \le \lVert\Lambda^n\rVert
  +(D-1)n^{D-1}\lVert\Lambda\rVert^{n-D+1}
    \max\{\lVert N\rVert,\lVert N\rVert^{D-1}\}.
\end{align}
The local source is
\path{Notes/WolfNoteTexSource/ch08_distance_measures.tex}, lines~1189--1210,
especially Equation~(8.103) and its derivation from Equation~(8.105).

\section{The exponent defect}

When $n<D-1$, the exponent $n-D+1$ is negative.  The source gives no
convention for this negative power.  In particular, if $\Lambda=0$, then the
factor $\lVert\Lambda\rVert^{n-D+1}$ is not defined in the usual real-power
interpretation.  Replacing the exponent by truncated natural subtraction
would silently change the displayed statement.

The derivation from Equation~(8.105) also factors
\begin{align}
  \lVert\Lambda\rVert^{n-k}
  =\lVert\Lambda\rVert^{n-(D-1)}
   \lVert\Lambda\rVert^{(D-1)-k},
\end{align}
which uses natural exponents only when $D-1\le n$ (and $k\le D-1$).
Thus the proof printed in the notes establishes the displayed natural-power
bound in that range, but does not supply an all-$n$ interpretation.

\section{Formalization scope}

Equation~(8.105), whose exponents satisfy $k\le n$, is formalized for every
$n\in\mathbb N$.  Equation~(8.103) is formalized with the explicit hypothesis
$D-1\le n$.  Its refined constant $(D-1)\binom{n}{D-1}$ is formalized under
the source's stated hypothesis $2(D-1)\le n$, which implies the first range
condition.  Both results retain the arbitrary submultiplicative seminorm
parameter; no particular matrix norm replaces it.  The corresponding Lean
declarations are:
\begin{itemize}
  \item \leanid{Matrix.wolf_eq_105_seminorm} for the unrestricted
    Equation~(8.105) word bound;
  \item \leanid{Matrix.wolf_eq_103} for the natural-power form under
    $D-1\le n$;
  \item \leanid{Matrix.wolf_eq_103_refined} for the refined constant under
    $2(D-1)\le n$.
\end{itemize}

Equation~(8.111) applies Equation~(8.103) to the Schur form of
$\widehat T-\widehat T_\varphi$, with matrix dimension $D=d^2$ and
$\lVert\Lambda\rVert_\infty=\mu$.  It therefore inherits the same restriction
$d^2-1\le n$ in the natural-power interpretation.  The corresponding coarse
and refined Schur-form estimates are
\leanid{Matrix.wolf_eq_111_schur_form} and
\leanid{Matrix.wolf_eq_111_schur_form_refined}.  These declarations establish
the estimate once the Schur decomposition and its spectral identification are
given; constructing that decomposition for $\widehat T-\widehat T_\varphi$
remains part of the channel-level theorem.

The specializations
\leanid{Matrix.wolf_eq_111_schur_form_subperipheralModulus} and
\leanid{Matrix.wolf_eq_111_schur_form_subperipheralModulus_refined} now use the
shared source-shaped $\mu$ from
\leanid{Module.End.subperipheralModulus} on the same $D$-dimensional coordinate
space as the Schur matrices.  They discharge $\mu\le1$, including the empty
subperipheral-spectrum convention.  They intentionally retain the hypothesis
$\lVert\Lambda\rVert_\infty=\mu$: constructing the unitary Schur representation
of $\widehat T-\widehat T_\varphi$ and proving this spectral identification
remain separate channel-level obligations.

\section{Verdict}

\textbf{Verdict: source scope restriction.}

This is not a missing Lean proof.  The
formal statement records exactly the range in which the printed
natural-exponent derivation is meaningful.  Eliminating the restriction would
require Wolf to specify an interpretation for negative powers, including the
case $\lVert\Lambda\rVert=0$.  Tracking: \href{https://github.com/LionSR/TNLean/issues/5837}{TNLean issue \#5837}.

\bibliographystyle{alpha}
\bibliography{references}

\end{document}
